Multimorbidity, defined as the co-occurrence of two or more clinical conditions, has emerged as a critical challenge in modern healthcare systems. With aging populations and improved survival rates, its prevalence continues to rise, affecting approximately 25\% of adults and over 65\% of those aged 65 years or older~\citet{barnett2012epidemiology}. Cardiovascular and cardiometabolic diseases represent a particularly significant burden, accounting for nearly 18 million deaths annually worldwide and frequently presenting alongside metabolic, renal, and systemic disorders~\citet{roth2020global}. Understanding the complex patterns of disease co-occurrence and their underlying risk factors is essential for developing targeted prevention strategies, optimizing treatment protocols, and improving patient outcomes in complex clinical scenarios~(\citealp{PradosTorres2014, Yan2022}). %Moreover, characterizing these relationships can reveal shared underlying biological mechanisms and identify potential therapeutic targets that extend across multiple conditions. 

The study of comorbidity networks has gained substantial attention as a powerful framework for characterizing disease relationships~\citet{DelValle2019}. These networks represent diseases as nodes and their associations as edges, providing a systems-level view of multimorbidity patterns~\citet{Barabasi2011}. Traditional approaches to comorbidity network construction typically rely on cross-sectional prevalence or incidence data, where disease associations are computed as simultaneous presence rates within the study cohort at the specific time point~(\citealp{hidalgo2009dynamic, goh2007human}). Recent methodological advances have incorporated electronic health records and administrative databases to construct large-scale disease networks, leveraging statistical measures such as relative risk, odds ratios, phi-correlation, or Pearson correlation to quantify disease relationships~(\citealp{jensen2014temporal, siggaard2020disease, VargasFernandez2026Network, Burke2023}). While these methods have identified clinically meaningful disease clusters and progression pathways, they face three critical limitations.

First, most existing methods discard temporal information by using aggregate prevalence statistics or single time-point measurements~\citet{jeong2017network, dervic2025comorbidity}. Although some approaches incorporate longitudinal data for directed network construction, they typically rely on simple pairwise co-occurrence counts rather than learning representations that jointly capture temporal progression patterns and semantic disease relationships. This results in incomplete characterizations that fail to reflect the dynamic nature of multimorbidity development. Second, disease associations estimated solely from co-occurrence patterns are vulnerable to confounding by shared risk factors~(\citealp{vanderweele2013definition, pearl2009causality}). Two diseases may frequently co-occur not due to direct pathophysiological connections, but because they share common antecedents such as lifestyle factors or other comorbid conditions. Traditional network construction methods that ignore this confounding produce biased association estimates, artificially inflating similarities between diseases with shared risk profiles. Third, typical pair-wise estimation methods produce fully connected weighted networks that require arbitrary global thresholding for sparsification~\citet{masuda2025}. The resulting pruned networks depend critically on threshold selection, potentially removing clinically relevant connections while retaining spurious ones. Although statistical "backboning" methods have been proposed~(\citealp{serrano2009extracting, foti2011nonparametric}), these approaches assume heterogeneous weight distributions and often fail to extract meaningful structure from networks with relatively homogeneous similarity values.

%While graphical models can identify conditional dependencies, their application to disease network inference has not been combined with explicit risk factor adjustment in temporal sequence contexts.



To address these limitations, we develop a comprehensive statistical framework for disease network inference from longitudinal patient data that integrates three methodological innovations. First, we leverage complete longitudinal disease sequences as input, capturing the chronological progression of disease manifestations over the life course. We adapt Word2Vec embeddings, originally developed for natural language processing~\citet{mikolov2013efficient}, to longitudinal clinical sequences. We fine-tuned pretrained biomedical embeddings~\citet{zhang2019biowordvec} to learn dense vector representations that jointly encode temporal co-occurrence patterns and semantic disease relationships. Second, we introduce an approach to confounding adjustment by jointly embedding diseases and risk factors within a single representation space. This is motivated by the distributional principle underlying Word2Vec-style models, where incorporating additional relevant context reshapes the local co-occurrence neighborhoods used for training, leading to richer and more informative semantic representations. We explicitly evaluated this reshaping by comparing risk-adjusted and unadjusted similarity networks to quantify confounding effects. We validate confounding removal by demonstrating that disease pairs with larger similarity changes systematically share more risk factors, confirming that observed decreases reflect genuine confounding adjustment rather than noise. Third, we employ Gaussian Graphical Models (GGM) with Lasso regularization~\citet{meinshausen2006high} to achieve network sparsification by estimating conditional independence structure, identifying direct disease associations while removing indirect relationships mediated through other diseases or risk factors. This conditional independence framework provides a statistically rigorous alternative to arbitrary thresholding, naturally producing sparse networks suitable for topological analysis and identification of hub diseases, improving clinical interpretability.

We apply this framework to construct a cardiometabolic comorbidities disease network using longitudinal data from UK Biobank, incorporating 24 primary cardiometabolic diseases and 76 risk factors. Our results demonstrate substantial confounding, with 91.3\% of disease pairs showing decreased association after risk adjustment. The sparse GGM-derived network preserves established pathophysiological relationships while revealing hub diseases with both high epidemiological burden and critical bridging functions between distinct biological domains. 

The remainder of this paper is structured as follows. Section~\ref{sec:data} describes the UK Biobank cohort, phenotype definitions, and data preprocessing. Section~\ref{sec:method} presents our methodological framework, including sequence construction, word embedding strategy, similarity-based network construction, confounding evaluation, and GGM-based sparsification. Section~\ref{sec:results} reports results for cardiometabolic disease networks, including confounding quantification and topological analysis. Section~\ref{sec:discussion} discusses methodological implications, limitations, and future directions.